\documentclass[a4paper,12pt]{article}
\usepackage[utf8]{inputenc}
\usepackage[T1]{fontenc}
\usepackage{lmodern}
\usepackage{graphicx}
\graphicspath{ {Images/} }
\usepackage{placeins}
\usepackage{epstopdf}
\usepackage{subfigure}
\usepackage{pstricks}
\usepackage{epsfig} % extra package for .eps graphics usage
\usepackage{pst-grad} % For gradients
\usepackage{pst-plot} % For axes
\usepackage{amsmath} % extra math symbols There are a lot of ams packages
\usepackage{fancyhdr} % to use fancy headers and footers --- requires for most professional docs
\usepackage{watermark} % to include a figure in the background - the front page for example
\usepackage{enumerate} % to enumerate your items in a list with any style
\usepackage{booktabs}
\usepackage{siunitx}
\usepackage{wallpaper}
\usepackage{afterpage}
\usepackage{color}
\definecolor{mygreen}{RGB}{28,172,0}
\definecolor{mylilas}{RGB}{170,55,241}
\usepackage{listings}
\lstset{language=Matlab,%
    %basicstyle=\color{red},
    breaklines=true,%
    morekeywords={matlab2tikz},
    keywordstyle=\color{blue},%
    morekeywords=[2]{1}, keywordstyle=[2]{\color{black}},
    identifierstyle=\color{black},%
    stringstyle=\color{mylilas},
    commentstyle=\color{mygreen},%
    showstringspaces=false,%without this there will be a symbol in the places where there is a space
    numbers=left,%
    numberstyle={\tiny \color{black}},% size of the numbers
    numbersep=9pt, % this defines how far the numbers are from the text
    emph=[1]{for,end,break},emphstyle=[1]\color{red}, %some words to emphasise
    %emph=[2]{word1,word2}, emphstyle=[2]{style},    
}
\numberwithin{equation}{section}
\numberwithin{figure}{section}
\numberwithin{table}{section}
\usepackage[toc,page,title,titletoc]{appendix}
\usepackage{amsmath}
\usepackage[
backend= bibtex,
style=ieee,
sorting=none
]{biblatex}
\addbibresource{references.bib}

\usepackage[colorlinks=true, linkcolor=blue, citecolor=red, urlcolor=blue]{hyperref}
\parindent=0pt
\parskip=6pt

\renewcommand{\contentsname}{Table of Contents}
%-------------------------------------------------------------------------------
%-------------------------------------------------------------------------------
%-------------------------------------------------------------------------------

\title{\vspace{3 cm}
\includegraphics[scale=0.5]{fronttext}\vspace*{2cm} \\
\Huge{\textsc{
    Milestone 2 \\
    Control Theory 2   } \\ \Large
    EERI 418 }
}
		

\author{by:\\ \begin{tabular}{l  r}
    W Bisschoff  & 26217406
\end{tabular}	\\
\\
\textbf{North-West University  Potchefstroom Campus}\\ \\
\begin{Large} 
\begin{tabular}{l l} \\
    Lecturer: & Dr Kenny Uren 
\end{tabular} 
\end{Large} \\
\vspace*{0.8cm}
}

%-------------------------------------------------------------------------------
%-------------------pagestyle stuff & headers & footers ------------------------
%-------------------------------------------------------------------------------
\pagestyle{fancy} % other pagestyles can be used
%\headsep=0.4cm
\headheight=15pt % Have enough space so that the header can fit without warnings
\headwidth=6.7in % the width of the header 
\textwidth=6.6in % the width of the text
\oddsidemargin=0in % the margin indent on odd pagenumbers
\evensidemargin=0in % the margin indent on even pagenumbers

\rhead{\begin{pspicture}(0,0)(0,0)
			\includegraphics[scale=0.25]{logo_fade.png}
		\end{pspicture}
	}
\chead{} % center header
\lhead{\textbf{\textsf{\footnotesize FACULTY OF ENGINEERING}}} % left header
%-------------------------------------------------------------------------------
%-------------------------------------------------------------------------------
%-------------------------------------------------------------------------------


\lfoot{Milestone 2} 


%-------------------------------------------------------------------------------
%-------------------------------------------------------------------------------
%-------------------------------------------------------------------------------
\cfoot{\today} %The current date will be put there
\rfoot{\thepage} %pagenumber
\renewcommand{\headrulewidth}{1pt} % to increase the header or footer line size
\renewcommand{\footrulewidth}{1pt}
%-------------------------------------------------------------------------------
%-------------------------------------------------------------------------------
%-------------------------------------------------------------------------------



%*********************************************************************************
%----------------------------Beginning of the document----------------------------
%*********************************************************************************
\begin{document}
%\sffamily   % include this for sans serif font family
%\fontfamily{pcr} more ways of changing font see: http://tex.loria.fr/general/new/fntguide.html

%\numberwithin{equation}{subsection}
\pagenumbering{roman} % first pages are number in Roman numerals

\ThisULCornerWallPaper{1}{frontcorner.eps}
\maketitle % to display the author and title that was defined earlier

\thispagestyle{empty} % this page should not have header or footers
\pagebreak % break the page here
%-------------------------------------------------------------------------------
%-------------------------------------------------------------------------------
%-------------------------------------------------------------------------------

\section*{Abstract}
Abstract
\pagebreak


\tableofcontents

\newpage{}
{
\let\oldnumberline\numberline%
\renewcommand{\numberline}{\figurename~\oldnumberline}%
\listoffigures
}
\newpage{}
{
\let\oldnumberline\numberline%
\renewcommand{\numberline}{\tablename~\oldnumberline}%
\listoftables
}



% The beginning of the sections
\newpage
\pagenumbering{arabic}

\section{Introduction}
Introduction


\clearpage

\section{Introduction}\label{introduction}

A DC motor is to be characterized and a mathematical model is to be
derived.

\section{The Method}\label{the-method}

Calculate the armature resistance $R$ by measuring the armature
current $I$ when supplying voltage $V$ to a short-circuit and
measuring the armature voltage $E$ when supplying voltage to the open
circuit.

Applying an AC voltage and determining the armature voltage and current
allows for the calculation of the armature impedance $Z$. The
inductance $L$ can be calculated from the reactance $X$ of the
impedance.

The back-EMF constant $K_e$ can be determined from

 \begin{equation} 
K_e = \frac{\Delta E}{\Delta\omega}, 
 \end{equation} 
 and the torque constant
$K_t$ can be determined as $K_t = K_e$.

The damping coefficient $b$ can be determined from

 \begin{equation} 
b\omega = K_tI_a, 
 \end{equation} 
 which is derived from 
 \begin{equation} 
T=K_ti
 \end{equation} 
 and

 \begin{equation} 
T = J\dot\omega +b\omega, 
 \end{equation} 
 while keeping the system at a steady
state so that $\dot\omega = 0$.

The motor inertia $J$ can be calculated from

 \begin{equation} 
J = -\frac{T}{\frac{dw}{dt}}.
 \end{equation} 


\section{Results}\label{results}

The resistance $R = 10,42\Omega$ can be calculated from Table 1.

\begin{longtable}[]{@{}lll@{}}
\toprule
$V_f$ & $E_a$ & $I_a$\tabularnewline
\midrule
\endhead
0 & 19.10 & 1,90\tabularnewline
4 & 23,03 & 2,25\tabularnewline
8 & 26,19 & 2,58\tabularnewline
12 & 29,82 & 2,81\tabularnewline
16 & 34,16 & 3,13\tabularnewline
20 & 37,67 & 3,48\tabularnewline
\bottomrule
\end{longtable}

The back-EMF constant and the torque constants can be calculated as
$K_t = K_e = 0.728$, the damping coefficient can be calculated as
$b = 0.014$. The inductance $L = 0.15$.

\begin{longtable}[]{@{}lll@{}}
\toprule
Ea & Ia & w\tabularnewline
\midrule
\endhead
24 & 1,186 & 12,15\tabularnewline
38 & 1,408 & 28,48\tabularnewline
55 & 1,605 & 49,95\tabularnewline
67 & 1,67 & 66,71\tabularnewline
76 & 1,705 & 79,27\tabularnewline
84 & 1,72 & 89,12\tabularnewline
92 & 1,738 & 100,16\tabularnewline
96 & 1,747 & 109,01\tabularnewline
105 & 1,752 & 120,32\tabularnewline
115 & 1,78 & 135,40\tabularnewline
128 & 1,76 & 155,72\tabularnewline
\bottomrule
\end{longtable}

The inertia $J = 0,022$.

The transfer function $P(s)$ can be determined as


 \begin{equation} \label{eq:tf}
P(s)=\frac{\omega(s)}{V(s)}=\frac{K_t}{(Js+b)(Ls+R)+K_eK_t}.

 \end{equation} 

 The step response of the transfer function in equation~\eqref{eq:tf} can be seen in Figure~\ref{fig:step}.

 \begin{figure}
    \centering
    \includegraphics{stepplot.eps}
    \caption{Step Response}
    \label{fig:step}
 \end{figure}




\end{document}